\documentclass[conference]{IEEEtran}
\usepackage[spanish]{babel}
\usepackage[utf8]{inputenc}
\usepackage[T1]{fontenc}
\usepackage{amsmath}
\usepackage{booktabs}
\usepackage{url}

\title{Notacion de Landau para Analisis Asintotico}
\author{\IEEEauthorblockN{Ronquillo Nunez Braulio}
\IEEEauthorblockA{Analisis y Diseno de Algoritmos}}

\begin{document}
\maketitle

\begin{abstract}
La notacion de Landau es una herramienta central para describir el crecimiento
de funciones asociadas al costo de un algoritmo. En lugar de estudiar tiempos
medidos en una computadora especifica, el analisis asintotico observa el
comportamiento cuando el tamano de entrada crece. Este reporte presenta las
definiciones de $O$, $\Omega$ y $\Theta$, su interpretacion con limites y la
forma en que se implemento un programa para comparar polinomios mediante sus
terminos dominantes.
\end{abstract}

\begin{IEEEkeywords}
Landau, Big-O, complejidad, analisis asintotico, polinomios.
\end{IEEEkeywords}

\section{Introduccion}
El analisis de algoritmos busca abstraer detalles de plataforma como lenguaje,
procesador o memoria cache. Para lograrlo, se modela el costo de un algoritmo
como una funcion $T(n)$, donde $n$ representa el tamano de la entrada. La
notacion de Landau permite expresar cotas de crecimiento sin depender de
constantes multiplicativas ni de terminos que desaparecen frente al termino
dominante.

Segun NIST, $f(n)=O(g(n))$ indica que $f$ queda acotada por un multiplo
constante de $g$ a partir de cierto punto \cite{nist_big_o}. MathWorld resume
los simbolos de Landau como una familia de notaciones para describir funciones
cuando la variable tiende a infinito \cite{mathworld_landau}. En algoritmos,
esto se traduce en una comparacion teorica entre soluciones.

\section{Marco Teorico}
Sea $f(n)$ el costo de un algoritmo y $g(n)$ una funcion de referencia. La
cota superior asintotica se define como
\begin{equation}
f(n)=O(g(n)) \iff \exists c>0,\exists n_0:0\leq f(n)\leq c g(n)
\end{equation}
para todo $n\geq n_0$. Esta relacion no afirma igualdad exacta; solo indica
que $g$ domina a $f$ hasta una constante. Por otro lado, $\Omega(g(n))$ se
interpreta como cota inferior y $\Theta(g(n))$ como cota ajustada. NIST define
$\Theta$ mediante constantes $c_1$, $c_2$ y $k$ tales que
\begin{equation}
0\leq c_1g(n)\leq f(n)\leq c_2g(n)
\end{equation}
para todo $n\geq k$ \cite{nist_theta}.

En el caso de polinomios, el termino de mayor exponente determina la clase
asintotica. Por ejemplo,
\begin{equation}
4n^3+2n^2+7n+9=\Theta(n^3).
\end{equation}
Los terminos menores no cambian la clase porque su cociente contra $n^3$
tiende a cero.

\section{Criterio por Limites}
Una forma practica de comparar dos funciones positivas es analizar
\begin{equation}
L=\lim_{n\to\infty}\frac{f(n)}{g(n)}.
\end{equation}
Si $L=0$, entonces $f$ crece menos que $g$. Si $0<L<\infty$, ambas funciones
son del mismo orden. Si $L=\infty$, entonces $f$ domina a $g$. Este criterio
es especialmente claro en polinomios, exponenciales y logaritmos.

\begin{table}[htbp]
\caption{Interpretacion del cociente asintotico}
\centering
\begin{tabular}{ll}
\toprule
Resultado del limite & Relacion \\
\midrule
$0$ & $f(n)=O(g(n))$ estricta \\
$c>0$ & $f(n)=\Theta(g(n))$ \\
$\infty$ & $g(n)=O(f(n))$ estricta \\
\bottomrule
\end{tabular}
\end{table}

\section{Implementacion}
El programa implementado recibe dos polinomios. Cada polinomio se representa
como una lista de pares $(coeficiente, exponente)$. El algoritmo localiza el
termino dominante de cada funcion y compara primero los exponentes. Si los
exponentes son distintos, el polinomio con mayor exponente domina. Si son
iguales, se compara el cociente de coeficientes y se concluye una relacion
$\Theta$.

\section{Analisis de Complejidad}
Si cada polinomio contiene $m$ terminos, basta recorrer las listas para
encontrar el mayor exponente. Por tanto, el tiempo es $O(m)$ y el espacio
adicional es $O(1)$. La decision es correcta para polinomios porque el termino
dominante controla el limite del cociente.

\section{Conclusion}
La notacion de Landau formaliza la comparacion entre algoritmos cuando el
tamano de entrada crece. En esta practica, el caso de polinomios permite ver
la idea esencial: constantes y terminos menores desaparecen frente al termino
dominante. Por eso, un programa que compara exponentes principales reproduce
el criterio teorico de crecimiento asintotico.

\begin{thebibliography}{00}
\bibitem{nist_big_o} P. E. Black, ``big-O notation,'' Dictionary of Algorithms and Data Structures, NIST, 2019. [Online]. Available: \url{https://www.nist.gov/dads/HTML/bigOnotation.html}
\bibitem{nist_theta} P. E. Black, ``Theta,'' Dictionary of Algorithms and Data Structures, NIST, 2016. [Online]. Available: \url{https://www.nist.gov/dads/HTML/theta.html}
\bibitem{mathworld_landau} E. W. Weisstein, ``Landau Symbols,'' MathWorld--A Wolfram Resource. [Online]. Available: \url{https://mathworld.wolfram.com/LandauSymbols.html}
\bibitem{mit_algorithms} E. Demaine, J. Ku, and J. Solomon, ``6.006 Introduction to Algorithms: Lecture Notes,'' MIT OpenCourseWare, 2020. [Online]. Available: \url{https://ocw.mit.edu/courses/6-006-introduction-to-algorithms-spring-2020/pages/lecture-notes/}
\end{thebibliography}

\end{document}
